\documentclass[11pt]{article}
\usepackage{../homework}

\profname{Dr. Stanislaw Szarek}
\classcode{MATH 423 Measure Theory}
\assignnum{Assignment 10}
\duedate{09 November 2026}

\begin{document}
\begin{problem}{2.4.34(a,b)}
  Suppose \( |f_n| \leq g \in L^1 \) and \( f_n \rightarrow f \) in measure.
  \begin{enumerate}[label=(\alph*)]
    \item \( \int f = \lim \int f_n \).
          \begin{proof}
            Since \( f_n \rightarrow f \) in measure, there exists, by Theorem 2.30, a  subsequence \( f_{n_k} \) such that \( f_{n_k} \rightarrow f \) a.e.
            This in turn gives that
            \begin{equation}
            \label{eq:1}
            \int f_{n_k} \rightarrow \int f,
          \end{equation}
          but this is not quite what we want to prove.
          We instead do the following.
          First of all, since \( f_n \rightarrow f \) in measure, for every subsequence \( f_{n_k} \), \( f_{n_k} \rightarrow f \) in measure.
          This is essentially by definition, since if for every \( \epsilon > 0 \),
          \begin{equation}
          \label{eq:2}
          \mu(\{x:|f_n(x)-f(x)| \geq \epsilon\}) \xrightarrow{n \rightarrow \infty} 0,
        \end{equation}
        then \( \mu(\{x: |f_{n_k} - f| \geq \epsilon\}) \) is a subsequence of \( \mu(\{x: |f_n - f| \geq \epsilon\}) \) in \( \mathbb{R} \), so that \( \mu(\{x: |f_{n_k} - f| \geq \epsilon\}) \rightarrow 0 \).
      From this, we apply Theorem 2.30 to \( f_{n_k} \) to obtain an additional subsequence \( f_{n_{k_i}} \rightarrow f \) a.e.
      Phrasing this in terms of sequences of real numbers, we have
      \begin{equation}
      \label{eq:4}
      I_n = \int f_n, \quad \text{ for every } I_{n_k} = \int f_{n_k}, \ \exists I_{n_{k_i}} \rightarrow I = \int f
    \end{equation}
    where the last convergence comes from \( f_{n_{k_i}} \rightarrow f \) a.e. implying \( \int f_{n_{k_i}} \rightarrow \int f \) (since \( f_n \) is majorized by \( g \) so in particular any subsequence is).
    We now prove the following lemma:
    \begin{lem}
      \label{lem:1}
      Let \( (a_n) \) be a sequence in \( \mathbb{R} \) such that for every subsequence \( (a_{n_k}) \) there exists \( a_{n_{k_i}} \rightarrow l \).
      Then \( a_n \rightarrow l \).
    \end{lem}
    \begin{proof}
      Fix some \(\epsilon > 0\), and set \( A = \{a_n: |a_{n} - l| \} \geq \epsilon\).
      If \( A \) is a finite set, we are done, all but finitely many \( a_n \) are \(\epsilon\)-close to \( l \), that is, \( a_n \rightarrow l \).
      If, on the other hand, \( A \) is finite, it defines a subsequence \( a_{n_k} \) of \( a_n \).
      Then, by assumption, there is a further subsequence \( a_{n_{k_i}} \rightarrow l \).
      That is, infinitely many elements of \( A \) are actually \(\epsilon\)-close, a contradiction.
    \end{proof}

    Since our sequences \( I_n, I_{n_k}, I_{n_{k_i}} \) satisfy the assumptions of the lemma with \( l = \int f \), it is immediate that
    \begin{equation}
    \label{eq:5}
    \int f_n \rightarrow f,
  \end{equation}
  as desired.
  \end{proof}
    \item \( f_n \rightarrow f \) in \( L^1 \).
          \begin{proof}
            We want to show that \( \int |f_n - f| \rightarrow 0 \).
            We proceed quite similarly to part (a).
            Part (a) showed that \( f \in L^1(\mu) \) since \( \int f = \lim \int f_n < \infty \).
            Again, for any subsequence \( f_{n_k} \), \( f_{n_k} \rightarrow f \) in measure, from which we can find a further subsequence \( f_{n_{k_i}} \rightarrow f \) a.e.
            This in particular means that
            \begin{equation}
            \label{eq:21}
            |f_{n_{k_i}} - f| \rightarrow 0,
          \end{equation}
          and that, by Exercise 2.2.15, since \( |f_{n_{k_i}} - f| \in L^+ \) and decreases to \( 0 \),
          \begin{equation}
            I_{n_{k_i}} \coloneq \int |f_{n_{k_i}} - f| \rightarrow 0.
          \end{equation}
          This again fulfills the assumptions of \cref{lem:1}, so that
          \begin{equation}
          \label{eq:22}
          I_n = \int |f_n - f| \rightarrow 0,
        \end{equation}
        as desired.
          \end{proof}
  \end{enumerate}
\end{problem}

\begin{problem}{2.5.46}
  Let \( X = Y = [0,1] \), \( \mathcal{M} = \mathcal{N} = \mathcal{B}_{[0,1]} \), \( \mu = m \), \( \nu = \text{counting measure} \).
  If \( D = \{(x,x) : x \in [0,1]\} \) is the diagonal in \( X \times Y \), then \( \int \int \chi_D \dd{\mu} \dd{\nu} \), \( \int \int \chi_D \dd{\nu}\dd{\mu} \), and \( \int \chi_D \dd{(\mu \times \nu)} \) are all unequal.
\end{problem}
\begin{proof}
  To begin with, let \( \bigcup_{k=1}^{\infty} E_k \times F_{k} \) be a cover of \( D \).
  Since \( D \) is compact, we can take a finite subcover \( D \subset \bigcup_{k=1}^n E_k \times F_k \).
  Now, if \( (a,b) \in D \), then \( (a,b) \in \bigcup_{k=1}^{n} E_k \times F_k \), but by definition of \( D \), \( a = b = x \), so that \( (x,x) \in \bigcup_{k=1}^{n} E_k \times F_k \), meaning \( x \in E_k \cap F_k \neq \emptyset \) and \( D \subset \bigcup_{k=1}^{n} (E_k \cap F_k) \times (E_k \cap F_k) \).

  Now, this means that for every \( x \in [0,1] \), \( x \in \bigcup_{k=1}^{n} E_k \cap F_k \), or in other words,
  \begin{equation}
  \label{eq:6}
  [0,1] \subset \bigcup_{k=1}^{n} E_k \cap F_k,
\end{equation}
which implies
\begin{equation}
\label{eq:7}
m([0,1]) = 1 \leq m(\bigcup_{k=1}^{n} E_k \cap F_k).
\end{equation}
Since every countable set has Lebesgue measure \( 0 \), that means at least one pair \( E_k \cap F_k \) is uncountable.
Furthermore,
\begin{equation}
\label{eq:11}
(\mu \times \nu)\left( \bigcup_{k=1}^n (E_k \cap F_k) \times (E_k \cap F_k) \right) = \sum\limits_1^n \mu(E_k \cap F_k)\nu(E_k \cap F_k) = \infty,
\end{equation}
since the counting measure of any uncountable set is infinite. (There is an implicit disjointification argument here).

Now, by definition,
\begin{equation}
\label{eq:8}
\begin{aligned}
  \int \chi_D \dd{\mu \times \nu} &= (\mu \times \nu)(D) = \inf \left\{(\mu \times \nu)\left(\bigcup_1^{\infty} R_k \right):R_k\text{ rectangles covering D} \right\}.
\end{aligned}
\end{equation}
However, as we established, any cover of \( D \) by rectangles has measure \(\infty\), so that
\begin{equation}
\label{eq:12}
\int \chi_D \dd{\mu \times \nu} = \infty.
\end{equation}

The other two integrals are much easier.
\begin{equation}
\label{eq:9}
\begin{aligned}
  \iint \chi_D \dd{\mu} \dd{\nu} &= \int \left( \int \chi^{y}_D(x,y) \dd\mu(x) \right)\dd\nu(y) \\
                      &= \int \left( \mu(D^y) \right) \dd\nu(y) = \int \left( \mu(\{x: (x,y) \in D\}) \right) \dd{\nu} = \int \left( \mu(\{y\}) \right) \dd{\nu} \\
  &= \int 0 \dd{\nu} = 0.
\end{aligned}
\end{equation}

\begin{equation}
\label{eq:10}
\begin{aligned}
  \iint \chi_D \dd{\nu}\dd{\mu} &= \int \left( \int (\chi_{D})_x \dd\nu(x) \right) \dd\mu(y) = \int \left( \nu(D_x) \right) \dd\mu \\
                     &= \int \nu(\{x: (x,y) \in D\}) \dd\mu = \int \nu(\{y\}) \dd\mu = \int \dd{\mu} \\
  &= \mu([0,1]) = 1.
\end{aligned}
\end{equation}
\end{proof}

\begin{problem}{2.5.48}
  Let \( X = Y = \mathbb{N} \), \( \mathcal{M} = \mathcal{N} = \mathcal{P}(\mathbb{N}) \), \( \mu = \nu = \text{counting measure} \).
  Define \( f(m,n) = 1 \) if \( m = n \), \( f(m,n) = -1 \) if \( m = n+1 \), and \( f(m,n) = 0 \) otherwise.
  Then \( \int |f| \dd{\mu \times \nu} = \infty \), and \( \iint f \dd{\mu}\dd{\nu} \) and \( \iint f \dd{\nu}\dd{\mu} \) exist and are unequal.
\end{problem}
\begin{proof}
  Let \( A \in \mathbb{N} \times \mathbb{N} \) be an arbitrary set.
  We can write \( A = \bigsqcup E_i \times F_i \) for rectangles \( E_i \times F_i \).
  Suppose that \( E_i \) contains \( q \) points and \( F_i \) contains \( p \) points.
  Then \( \mu(E_i) = q \), \( \mu(F_i) = p \), and \( (\mu \times \nu)(E_i \times F_i) = qp \).
  Since \( E_i \times F_i \) is the rectangle generated by a ``base'' of \( q \) points and a ``height'' of \( p \) points, it has \( qp \) points, so \( \mu \times \nu \) is a counting measure. (This is essentially the statement that \( \text{card}(E \times F) = \text{card}(E)\text{card}(F) \).)

  For the two iterated integrals, consider the \( x \) and \( y \) sections of \( f \),
  \begin{equation}
  \label{eq:14}
  f_x(y) = \begin{cases}
    1 &\text{if } y = x\\
    -1 &\text{if } y + 1 = x \\
    0 &\text{else}
  \end{cases},
  \end{equation}
  \begin{equation}
  \label{eq:13}
  f^y(x) = \begin{cases}
    1 &\text{if } y = x \\
    -1 &\text{if } x - 1 = y \\
    0 &\text{else}
  \end{cases}.
\end{equation}
Then notice that if \( x = 1 \in \mathbb{N} \), \( x - 1 = 0 \not\in \mathbb{N} \).
Thus with \( x = 1 \),
\begin{equation}
\label{eq:17}
\sum\limits_{y=1}^{\infty} f_x(y) = 1.
\end{equation}
However, for all \( y \in \mathbb{N} \), \( y + 1 \in \mathbb{N} \), so that for every \( y \),
\begin{equation}
\label{eq:18}
\sum\limits_{x=1}^{\infty} f^y(x) = 1 - 1 = 0.
\end{equation}
These two equalities show that
\begin{equation}
\label{eq:19}
\begin{aligned}
  \iint f \dd{\mu}(x)\dd{\nu}(y) &= \int \left( \int f^y(x)\chi_{\{y\}} \dd{\mu}(x) \right) \dd{\nu}(y) \\
  &= \int \left( \sum\limits_{x=1}^{\infty} f^y(x)\chi_{\{y\}} \right) \dd{\nu}(y) = \int 0\dd{\nu} = 0,
\end{aligned}
\end{equation}
but
\begin{equation}
\label{eq:20}
\begin{aligned}
  \iint f \dd{\nu}(y)\dd{\mu}(x) &= \int \left( \int f_x(y)\chi_{\{x\}}\dd{\nu}(y) \right) \dd{\mu}(x) \\
  &= \int \left( \sum\limits_{y=1}^{\infty} f_x(y)\chi_{\{x\}} \right) \dd{\mu}(x) = \int \chi_{\{1\}} \dd{\mu}(x) = \mu(\{1\}) = 1.
\end{aligned}
\end{equation}
\end{proof}
\end{document}
